\documentclass{article}
\usepackage{mathnotes}

\title{Random Variables and Probability Distributions}
\description{Mathematical framework for discrete and continuous random variables with their probability mass functions and cumulative distribution functions.}

\begin{document}

\section{Random Variables and Probability Distributions}

\subsection{Random Variables}

\begin{definition}[Random Variable]\synonyms{"variate"}
	A \textbf{random variable} is a \@{function} that maps from a \@{sample-space} to a \@{state-space}. Often, but not always, the state space is a \@{set} of \@{real} numbers. For \@{sample-space} $\Omega$, we have $X: \Omega \to \mathcal{X}$ and for $\omega \in \Omega$, we have $x = X(\omega),$ where $x$ is a \@{value} of the random variable and an element of the \@{state-space} $\mathcal{X}$.

	In other words, a random variable is a function whose \@{domain} is a \@{sample-space} and whose \@{codomain} is a \@{state-space} (called an \@{alphabet} when the state space is \@{discrete}).
\end{definition}

\@{Random variables} don't have to map to all \@{reals}; for example, \textbf{Bernoulli random variables} map to either 0 or 1.

\begin{definition}[State Space]\synonyms{"alphabet"}
	The \@{set} of values a \@{random-variable} can take is called its \textbf{state space}.
\end{definition}
\begin{note}
	While the \@{state-space} $\mathcal{X}$ is the \@{codomain} of a random variable, the actual values $X(\Omega)$ that are mapped to by the random variable are the \@{image} of the random variable. Furthermore, we often only care about the places were $p(X = x \in \mathcal{X}) > 0$ is non-zero.
\end{note}

Returning to \@{sample spaces} briefly:

\begin{definition}[Discrete Sample Space]
If a \@{sample space} contains a finite number of possibilities, or an unending \@{sequence} with as many elements as there are whole numbers (\@[countably infinite]{countable}), it is called a \textbf{discrete sample space}.
\end{definition}

\begin{definition}[Continuous Sample Space]
If a \@{sample space} contains an infinite number of possibilities equal to the number of points on a line segment, it is called a \textbf{continuous sample space}.
\end{definition}

\begin{definition}[Discrete Random Variable]
	A \@{random-variable} is called a \textbf{discrete random variable} if there exists a \@{countable} \@{set} $S \subseteq \mathcal{X}$ such that

\[ P \left ( \omega \in \Omega : X(\omega) \in S \right ) = 1. \]

In this case its state space is often called an alphabet.
\end{definition}

\begin{definition}[Absolutely Continuous Random Variable]\synonyms{"continuous random variable"}
	For a \@{real} valued \@{random-variable}

	\[ X : \Omega \to \mathbb{R}, \]

	we call $X$ a \textbf{continuous random variable} if there exists a \@{function}

	\[ f_X : \mathbb{R} \to [0, \infty) \]

	such that for every \@{measurable} set $A \subseteq \mathbb{R},$

	\[ P \left ( \omega \in \Omega : X(\omega) \in A \right ) = \int_{A}{f_X(x) dx} \]

	where $f_X$ is the \@{probability-density-function} of $X.$

	TODO: this definition excludes some continuous distributions like the Cantor distribution and in general isn't quite there yet in terms of measure theory. I need to revisit this when I've learned a bit more of measure theory.
\end{definition}

\@{Discrete random variables} tend to represent counted data, while \@{continuous random variables} tend to represent measured data.

\subsection{Probability Distributions}

\subsubsection{Discrete Probability Distributions}


\begin{definition}[Probability Distribution]\synonyms{"probability mass function", "probability function"}
A \@{discrete random variable} takes each of its values with a certain \@{probability}. The \@{function} $f(x)$ that gives the probability of each value $x$ of a \@{random variable} $X$ occurring is called the \textbf{probability function,} \textbf{probability mass function,} or \textbf{probability distribution.} More formally, the set of ordered pairs $(x, f(x))$ is a probability distribution of the discrete random variable $X$ if, for each possible outcome $x$:

\begin{enumerate}
\item $f(x) \geq 0$
\item $\sum_{x}{f(x)} = 1$
\item $P(X = x) = f(x)$
\end{enumerate}
\end{definition}

\begin{definition}[Cumulative Distribution Function]
Sometimes we want to know the \@{probability} that a \@{random variable} will be less than or equal to some \@{real number} $x.$ If we let $F(x) = P(X \leq x)$ for all real $x$, we define $F(x)$ to be the \textbf{cumulative distribution function} of the random variable $X.$ More formally, the \textbf{cumulative distribution function} $F(x)$ of a \@{discrete random variable} $X$ with a \@{probability distribution} $f(x)$ is:

\[ F(x) = P(X \leq x) = \sum_{t \leq x}{f(t)}, \quad -\infty < x < \infty. \]
\end{definition}

\subsubsection{Continuous Probability Distributions}

A \@{continuous random variable} has a \@{probability} of $0$ of exactly assuming any particular value. That's because, in any \@{interval} of \@{real-numbers}, there are infinitely many points, and so the denominator for the proportion is infinite. For example, there are infinitely many points between $1.1$ and $1.2$, and again, infinitely many points between $1.2$ and $1.21$ and so on, regardless of the level of precision.

So, instead of focusing on \@{probabilities} of particular values, we focus on probabilities of values falling in various \@{intervals}. For example, instead of asking what the probability that someone weighs $150$ pounds is, we can ask what is the probability that they weigh between $149$ and $151$ pounds.

\begin{definition}[Probability Density Function]
The \@{function} $f(x)$ is a \textbf{probability density function} (pdf) for the \@{continuous random variable} $X$, defined on the \@{reals}, if

\begin{enumerate}
\item $f(x) \geq 0$ for all $x \in \reals$
\item $\int_{-\infty}^{\infty} f(x) dx = 1 $
\item $P(a < X < b) = \int_{a}^{b} f(x) dx$
\end{enumerate}
\end{definition}

In other words, if we integrate $f(x)$ from $a$ to $b$ we will get the \@{probability} of $X$ taking on a value between $a$ and $b$.

From here, the \@{cumulative distribution function} $F(x)$ of a \@{continuous random variable} $X$ with density function $f(x)$ is attained by simply integrating $f(x)$ from $-\infty$ to $x$:

\[ F(x) = P(X \leq x) = \int_{-\infty}^{x} f(t)dt, \quad -\infty < x < \infty. \]

\end{document}
